\begin{lemma}[Relative ordered-zero construction]
\label[lemma]{lem:scheme-abel-inverse}
Put
\[
 \operatorname{Ord}_{\mathrm{fac},\mathscr G}
 =\operatorname{Ord}_{\mathrm{fac}}
   \times_{\mathscr R_T^{(2)}}\mathscr G.
\]
Let
$\operatorname{Ord}_{\mathrm{div}}$ be the functor over $\mathscr G$ whose
$W$-points are ordered pairs of relative effective Cartier divisors
$(D_1,D_2)$ of degree three on $C_W/W$ such that
\[
 \operatorname{div}(s_W)=D_W+D_1+D_2
\]
and the unordered pair $\{\pi_*D_1,\pi_*D_2\}$ is the unordered pair of line
sections of $E_W$ defined by the two components of $Q_{s_W}$.  This functor is
represented by a finite \'etale double cover of $\mathscr G$.
The assignment $Q_s=\ell_1\ell_2\mapsto(D_1,D_2)$ is an isomorphism of finite
\'etale double covers
\[
 \operatorname{Ord}_{\mathrm{fac},\mathscr G}
 \xrightarrow{\sim}\operatorname{Ord}_{\mathrm{div}}.
\]
It commutes with arbitrary base change, including nonreduced base change.
Moreover, each $D_i$ is a relative effective Cartier divisor of degree three,
and $[\mathcal O(D_i)\otimes\kappa^{-1}]$ lies in $\cX_{\mathrm{sm}/\cB}$.
\end{lemma}

\begin{proof}
For the forward construction put
$T'=\operatorname{Ord}_{\mathrm{fac},\mathscr G}$ and write
$Q_s=\ell_1\ell_2$ in its universal ordering.  By definition of
$\mathscr G$, none of the tangency, branch, collision, multiple-zero, or Abel
speciality loci excluded in \cref{def:admissible-locus} occurs.  For the
remainder of the forward construction, every object with a subscript $T$ is
pulled back to $T'$; this base change is suppressed from the notation.
The line $\ell_i$ then
cuts out a relative finite \'etale divisor $B_i\subset E_T$ of degree three.
Moreover
\[
 U_i=C_T\times_{E_T}B_i\longrightarrow B_i
\]
is finite \'etale of degree three, because $B_i$ avoids the branch divisor.

The zero over $B_i$ is selected scheme-theoretically as follows.  \'Etale-locally on
$B_i$, split $U_i=B_i\sqcup B_i\sqcup B_i$ and trivialize the restriction
of $3\kappa$, using the induced trivialization of its norm line.  If
$a_1,a_2,a_3\in\mathcal O_{B_i}$ are the three values of $s$, compatibility
of the finite locally free norm with base change gives
\begin{equation}
 a_1a_2a_3=\Nm_\pi(s)|_{B_i}=0
\label{eq:split-norm-zero}
\end{equation}
as an equality in the structure sheaf, not merely at geometric points.
The collision exclusions imply that locally two values, say $a_2,a_3$, are
nonzero on every geometric fibre.  Their vanishing loci are therefore empty,
so they are units in the corresponding structure sheaf.  Equation
\eqref{eq:split-norm-zero} then forces $a_1=0$, also in nilpotent directions.
Consequently
\[
 Z_i=Z(s)\times_{E_T}B_i\subset U_i
\]
is locally exactly one open-and-closed sheet of $U_i/B_i$.  Thus
$Z_i\to B_i$ is finite \'etale of degree one and hence is an isomorphism.

Let $t$ be a local equation for $B_i$ in $E_T$.  Since $B_i$ occurs with
multiplicity one in the norm divisor and the other norm factors are
disjoint from it, $\Nm_\pi(s)=t u$ with $u$ a unit along $B_i$.  \'Etale-locally
on a neighbourhood of $B_i$ in the complement of the branch divisor, split
the cover $C_T\to E_T$.  On the selected sheet the other two values of $s$
remain units, so the same local product formula shows that the selected
value is $t$ times a unit.
The zero is therefore simple.  The image $D_i\subset C_T$ of $Z_i$ is
finite \'etale of degree three over $T'$; as a finite \'etale multisection of
the smooth relative curve $C_T/T'$, it is a relative effective Cartier
divisor.

The pairwise disjoint divisors $D,D_1,D_2$ lie in $Z(s)$ and have total
relative degree nine, equal to $\deg(3\kappa)$.  After dividing by their
Cartier equations, the residual section has degree zero and is nonvanishing
on every geometric fibre.  Nakayama's lemma makes it a trivialization, so
there is no vertical or nilpotent residual contribution and
\begin{equation}
 Z(s)=D+D_1+D_2
\label{eq:relative-zero-sum}
\end{equation}
as relative Cartier divisors.  Conversely, an ordered divisor datum
determines the unique projective section $[s]$ with zero divisor
$D+D_1+D_2$, and taking its norm recovers the ordered line sections $B_i$
and hence the ordered factors $\ell_i$.  Thus these assignments are inverse
already at the divisor level.

It remains to compare this divisor construction with the relative Abel
incidence.  Let
\[
 \alpha:C_T^{(3)}\longrightarrow\operatorname{Pic}^3(C_T/T)
\]
be the relative Abel map and let
$\mathcal F=\Nm^{-1}(\eta^3)\subset\operatorname{Pic}^3(C_T/T)$.  The norm fibre is
the cartesian base change
\[
 N_T=C_T^{(3)}\times_{\operatorname{Pic}^3(C_T/T)}\mathcal F.
\]
Let $W_3^{(1)}\subset\operatorname{Pic}^3(C_T/T)$ be the open on which $h^0=1$
and the pushforward of a Poincar\'e bundle is locally free of rank one and
commutes with base change.  The standard
projective-bundle description of the Abel map and cohomology and base change
then identify its source with the projectivization of that rank-one bundle.
Thus the map is an isomorphism and remains so after arbitrary base change.
Fibrewise, this is the $h^0=1$ case of Riemann--Kempf
\cite{Kempf1973}:
\[
 \alpha^{-1}(W_3^{(1)})\xrightarrow{\sim}W_3^{(1)}.
\]
Base changing this isomorphism by
$\mathcal F\to\operatorname{Pic}^3(C_T/T)$ and
translating by $\kappa^{-1}$ gives
\[
 N_T^{(1)}\xrightarrow{\sim}X_T^{(1)}.
\]
Here, by definition, $X_T^{(1)}$ is the open part of the theta section
represented by line bundles with $h^0=1$.
Restrict both sides to the inverse images of
$X_T^{(1)}\cap X_{T,\mathrm{sm}/T}$, as imposed above.  Thus the divisor
construction descends to the relative smooth Abel incidence.

Zero schemes, fibre products, finite locally free norms, and the cartesian
Abel construction all commute with arbitrary base change.  Hence both
composites preserve the universal zero schemes and are the identity as
morphisms.  The use of $C_T^{(3)}\simeq\operatorname{Hilb}^3(C_T/T)$ is only
representability; no normalization of the ordering cover or of pair fibres
is involved.
\end{proof}

Let $x_T:T\to\cM$ be the point represented by the universal divisor $D$,
so that $x_T=[\mathcal O(D)\otimes\kappa^{-1}]$.  Put
\[
 h=x-e_1w^2+e_2w-e_3.
\]
In the basis used for \eqref{eq:determinant-factorization}, direct
substitution gives
\begin{equation}
 s_0+s_1=h,\qquad s_2=wh.
\label{eq:structural-pencil}
\end{equation}
The divisor of $h$ is $D+D_{-x_T}$: it contains $D$, and its residual
degree-three divisor has Abel class $-x_T$ and is unique on the $h^0=1$
locus.

Define
\[
 \cI_T^\dagger=
 \{y\in X_T^{(1)}\cap X_{T,\mathrm{sm}/T}:
       -x_T-y\in X_T^{(1)}\cap X_{T,\mathrm{sm}/T}\}.
\]
The two relative Abel inverses give universal divisors $D_1,D_2$ on
$C_T\times_T\cI_T^\dagger$.  Their Abel classes satisfy
\[
 [\mathcal O(D+D_1+D_2)\otimes\kappa^{-3}]=0
 \quad\text{in }\operatorname{Pic}^0(C_T/T).
\]
Thus $\mathcal O(D_1+D_2)$ and $\mathcal O(3\kappa-D)$ differ only by a line
bundle pulled back from the base.  Their divisor sections determine a unique
projective relative section
\[
 \vartheta:\cI_T^\dagger\longrightarrow
 \PP_T\bigl(p_*\mathcal O_{C_T}(3\kappa-D)\bigr).
\]
The factorization through the residual octic holds scheme-theoretically.
Indeed, on the hyperplane $u=v$, \eqref{eq:structural-pencil} gives
\[
 s=h(u+zw),\qquad
 \operatorname{div}(s)=D+D_{-x_T}+D_\kappa
 \quad(D_\kappa\in|\kappa|).
\]
The two residual Abel classes are therefore $-x_T$ and $0$.  This cannot
occur on $\cI_T^\dagger$, where both residual classes lie in
$X_T^{(1)}\cap X_{T,\mathrm{sm}/T}$.  Hence the pullback of $u-v=0$ is
empty, and the pullback of $u-v$ is a unit in a local trivialization of the
tautological line bundle.  Taking norms gives $(u-v)F_8=0$, so $F_8=0$ on
$\cI_T^\dagger$, including nilpotent directions.  Thus $\vartheta$ factors
through $\mathscr R_T$.  Set
\[
 \cI_T^{\mathrm{adm}}=\vartheta^{-1}(\mathscr G).
\]

\begin{proposition}[Abel-incidence identification]
\label[proposition]{prop:abel-incidence}
There is an isomorphism over $\mathscr G$
\begin{equation}
 \Phi:\operatorname{Ord}_{\mathrm{fac},\mathscr G}
       \xrightarrow{\sim}\cI_T^{\mathrm{adm}},
 \qquad
 ([s],\ell_1,\ell_2)\longmapsto
 [\mathcal O(D_1)\otimes\kappa^{-1}].
\label{eq:abel-incidence-isomorphism}
\end{equation}
Interchanging $\ell_1$ and $\ell_2$ corresponds to the involution
$y\mapsto-x_T-y$.  The isomorphism commutes with arbitrary base change.
\end{proposition}

\begin{proof}
On $\operatorname{Ord}_{\mathrm{fac},\mathscr G}$, put
\[
 y_i=[\mathcal O(D_i)\otimes\kappa^{-1}]\in P_T.
\]
The divisor equality \eqref{eq:relative-zero-sum} gives
\begin{equation}
 x_T+y_1+y_2
  =[\mathcal O(D+D_1+D_2)\otimes\kappa^{-3}]=0.
\label{eq:abel-class-sum}
\end{equation}
Hence $y_2=-x_T-y_1$, and the ordered-zero construction of
\cref{lem:scheme-abel-inverse} defines the morphism $\Phi$.  The defining
conditions of $\mathscr G$ show that it lands in
$\cI_T^{\mathrm{adm}}$.
